\documentclass[11pt]{amsart}

\usepackage[T1]{fontenc}
\usepackage{lmodern}
\usepackage{microtype}
\usepackage{mathtools,amssymb,amsthm}
\usepackage{array}
\usepackage[hidelinks]{hyperref}

\hypersetup{
  pdftitle={A Seven-Vertex Bipartite Graph Whose Positive Matching Decomposition
            Number Lies Below Its Slope Invariant}
}

\newtheorem{theorem}{Theorem}
\newtheorem{lemma}[theorem]{Lemma}
\newtheorem{corollary}[theorem]{Corollary}
\theoremstyle{remark}
\newtheorem*{remark}{Remark}

\newcommand{\pmd}{\operatorname{pmd}}
\newcommand{\slope}{\mathrm{s}}
\newcommand{\Kinv}{\mathrm{K}}
\newcommand{\Gam}{\Gamma}

\title[A bipartite graph with $\pmd<\Kinv$]
{A Seven-Vertex Bipartite Graph Whose Positive Matching Decomposition Number
 Lies Below Its Slope Invariant}

\date{}

\subjclass[2020]{05C70, 13F55, 05C85}
\keywords{positive matching decomposition, bipartite graph, edge ideal, slope,
greedy algorithm}

\begin{document}

\begin{abstract}
Farrokhi D.~G., Gharakhloo and Yazdan Pour introduced the positive matching
decomposition number $\pmd(\Gam)$ of a graph and, for a bipartite graph, an
invariant $\Kinv(\Gam)$ obtained by minimising the output of a greedy
slope-based algorithm over all labellings of the two sides; they proved
$\Delta(\Gam)\le\pmd(\Gam)\le\Kinv(\Gam)$ and asked whether the second
inequality is ever strict. We answer this affirmatively with the theta graph
$\Gam=\Theta(2,2,4)$ on seven vertices: $\pmd(\Gam)=3$ by a decomposition
exhibited below and checkable by hand, while $\Kinv(\Gam)=4$ by running the
authors' own algorithm on all $3!\,4!=144$ labellings (all $288$, if the two
sides may also be exchanged).
\end{abstract}

\maketitle

\section{The question}

Let $\Gam$ be a finite simple graph. Following Farrokhi D.~G., Gharakhloo and
Yazdan Pour \cite{FGY}, a matching $M$ of $\Gam$ is \emph{positive} when the
corresponding sum of variables generates a prime ideal after inversion; we
never use that definition directly, only the combinatorial characterisation the
authors prove for it \cite[criterion (iii)]{FGY}, which we therefore take as
the working definition here:

\begin{equation}\label{eq:pos}
\parbox{0.86\textwidth}{a matching $M$ of $\Gam$ is \emph{positive} if and only
if for every non-empty $N\subseteq M$ the induced subgraph $\Gam[V(N)]$ on the
endpoints of $N$ has a vertex of degree $1$ whose incident edge lies in $N$.}
\end{equation}

The \emph{positive matching decomposition number} $\pmd(\Gam)$ is the least $p$
for which the edge set admits an \emph{ordered} partition
$E(\Gam)=M_1\uplus\dots\uplus M_p$ such that each $M_i$ is a positive matching
of $\Gam-(M_1\cup\dots\cup M_{i-1})$. Since a matching contains at most one
edge at each vertex and the parts cover $E(\Gam)$,
\begin{equation}\label{eq:delta}
  \Delta(\Gam)\le\pmd(\Gam).
\end{equation}

For a bipartite graph the authors introduce an upper bound through a greedy
algorithm driven by a numerical \emph{slope}. Fix a bipartition
$V(\Gam)=X\uplus Y$ and a labelling $\tau$, that is, an ordering
$X=\{x_1,\dots,x_m\}$ and $Y=\{y_1,\dots,y_n\}$, and put
$\slope(x_iy_j)=j-i$. Their Algorithm~1 produces parts $M_1,M_2,\dots$ from the
edge set; we quote it from \cite{FGY}, with the pseudocode macros set as words
and indentation.

\medskip
\begin{center}\small
\begin{tabular}{@{}l@{}}
\textsc{Require}: A labeled bipartite graph $\Gam$ with bipartition $X,Y$\\
\textsc{Ensure}: A positive matching decomposition of $\Gam$\\
$k\longleftarrow 0$\\
\textbf{while} $E(\Gam)\neq\varnothing$\\
\quad $k\longleftarrow k+1$\\
\quad $X'\longleftarrow X\setminus\{\text{isolated vertices of }X\}$\\
\quad $Y'\longleftarrow\varnothing$\\
\quad $M_k\longleftarrow\varnothing$\\
\quad \textbf{while} $X'\neq\varnothing$\\
\quad\quad $s\longleftarrow\min(\mathrm{Slopes}(\Gam[X'\cup Y]))$\\
\quad\quad $M\longleftarrow\{xy\in E(\Gam[X'\cup Y])\colon\ \mathrm{slope}(xy)=s\}$\\
\quad\quad $M_k\longleftarrow M_k\cup\{xy\in M\colon\ y\notin Y'\}$\\
\quad\quad $X'\longleftarrow X'\setminus\{x\colon\ xy\in M\}$\\
\quad\quad $Y'\longleftarrow Y'\cup\{y\colon\ xy\in M\}$\\
\quad \textbf{end while}\\
\quad $E(\Gam)\longleftarrow E(\Gam)\setminus M_k$\\
\textbf{end while}\\
\textbf{return} $M_1,\ldots,M_k$
\end{tabular}
\end{center}

\medskip
Two edges of equal slope $s$ have distinct $X$-endpoints and, since $j=s+i$,
distinct $Y$-endpoints, so the class $M$ selected at a pass of the inner loop is
itself a matching and the order in which its edges are examined is immaterial:
each part, and hence the number $\Kinv(\tau)$ of parts, is well defined. Each
part is a matching, since an $x$ leaves $X'$ as soon as it contributes and a
$Y$-endpoint is used at most once. Finally
\begin{equation}\label{eq:K}
  \Kinv(\Gam)=\min_\tau \Kinv(\tau),
\end{equation}
the minimum over all $m!\,n!$ labellings, and \cite{FGY} proves
$\pmd(\Gam)\le\Kinv(\Gam)$; with \eqref{eq:delta} this gives the chain
$\Delta\le\pmd\le\Kinv$ that precedes the question we answer. The source denotes
this invariant by a fraktur $\mathfrak{K}$; we use an upright $\Kinv$. The
question is \cite[Question~1]{FGY}:

\begin{quote}\sloppy
\textbf{Question 1.} Is there any bipartite graph $\Gam$
satisfying $\pmd(\Gam)<\Kinv(\Gam)$?
\end{quote}

\noindent Since $\Kinv$ is defined through the slope $j-i$, which is not symmetric in the
two sides, and since Algorithm~1 peels from $X$, the definition leaves open
whether a witness may also exchange $X$ and $Y$. Every value below is given for
both readings: $\Kinv$ minimised over the $m!\,n!$ labellings of one
orientation, and over all $2\,m!\,n!$ labellings when the sides may be
exchanged. Nothing in this paper depends on the choice.

\section{The witness}

Let $X=\{x_1,x_2,x_3\}$, $Y=\{y_1,y_2,y_3,y_4\}$ and
\begin{equation}\label{eq:gamma}
  E(\Gam)=\{\,x_1y_1,\;x_1y_2,\;\;x_2y_1,\;x_2y_3,\;x_2y_4,\;\;
              x_3y_2,\;x_3y_3,\;x_3y_4\,\},
\end{equation}
that is, the rows of the bipartite adjacency matrix are the bitmasks
$3,13,14$ over $y_1,\dots,y_4$. So $\Gam$ has $7$ vertices and $8$ edges,
degrees $2,3,3$ on $X$ and $2,2,2,2$ on $Y$, and $\Delta(\Gam)=3$. The two
vertices of degree~$3$ are $x_2$ and $x_3$, and deleting them leaves three
paths with $1$, $1$ and $3$ vertices, so
\[
  \Gam=\Theta(2,2,4),
\]
the graph obtained from two vertices by joining them with three internally
disjoint paths of lengths $2$, $2$ and $4$. Under the labelling
$0,1,2=x_1,x_2,x_3$ and $3,4,5,6=y_1,y_2,y_3,y_4$ its graph6 string is
\texttt{FEhb?}.

\begin{theorem}\label{thm:main}
For $\Gam=\Theta(2,2,4)$ as in \eqref{eq:gamma},
\[
  \pmd(\Gam)=3<4=\Kinv(\Gam),
\]
the value of $\Kinv$ being $4$ both over the $144$ labellings of the
orientation printed above and over all $288$ labellings when the two sides may
be exchanged. In particular Question~1 of \cite{FGY} has answer \emph{yes}.
\end{theorem}

The two halves are of different character and we treat them separately: the
value of $\pmd$ is settled by an exhibited decomposition and \eqref{eq:delta},
the value of $\Kinv$ by a finite exhaustion of the authors' own algorithm.

\subsection*{$\pmd(\Gam)=3$}

\begin{lemma}\label{lem:forest}
Let $M$ be a matching of a graph $\Gam$ such that $\Gam[V(M)]$ is acyclic.
Then $M$ is positive.
\end{lemma}

\begin{proof}
Let $\emptyset\neq N\subseteq M$ and $F=\Gam[V(N)]$. As $V(N)\subseteq V(M)$
and an induced subgraph of an induced subgraph is induced, $F$ is a subgraph of
$\Gam[V(M)]$ and hence a forest. Every edge of $N$ has both endpoints in
$V(N)$, so $N\subseteq E(F)$ and $F$ has an edge; a forest with an edge has a
vertex $\ell$ of degree~$1$. Now $\ell\in V(N)$ and $N$ is a matching, so
$\ell$ is an endpoint of exactly one edge $f\in N$, and $f\in E(F)$. Since
$\ell$ has degree $1$ in $F$, the unique edge of $F$ at $\ell$ is $f$. Thus
$F$ has a degree-$1$ vertex whose incident edge lies in $N$, which is
\eqref{eq:pos}.
\end{proof}

\begin{corollary}\label{cor:forest}
If $\Gam$ is a forest then every matching of $\Gam$ is positive and
$\pmd(\Gam)=\Delta(\Gam)$.
\end{corollary}

\begin{proof}
Induced subgraphs of a forest are forests, so Lemma~\ref{lem:forest} applies to
every matching. A forest is $\Delta$-edge-colourable, and each colour class is
a matching, hence a positive one; so $\pmd\le\Delta$, and \eqref{eq:delta}
gives equality.
\end{proof}

\begin{proof}[Proof that $\pmd(\Gam)=3$]
Put
\begin{align*}
  M_1&=\{x_2y_1,\;x_3y_3\},\\
  M_2&=\{x_1y_1,\;x_2y_4,\;x_3y_2\},\\
  M_3&=\{x_1y_2,\;x_2y_3,\;x_3y_4\},
\end{align*}
an ordered partition of the eight edges \eqref{eq:gamma} into three matchings.
The subgraph induced by $V(M_1)=\{x_2,y_1,x_3,y_3\}$ is the path
$y_1x_2y_3x_3$, which is acyclic, so $M_1$ is positive by
Lemma~\ref{lem:forest}. Deleting $M_1$ leaves exactly the path
\[
  \Gam-M_1 \;=\; y_1\,x_1\,y_2\,x_3\,y_4\,x_2\,y_3
\]
on all seven vertices, a tree with $\Delta=2$; by Corollary~\ref{cor:forest}
its edge set splits into two positive matchings, namely its two alternating
classes, which are $M_2$ and $M_3$. Hence $\pmd(\Gam)\le3$. Conversely
$\pmd(\Gam)\ge\Delta(\Gam)=3$ by \eqref{eq:delta}, so $\pmd(\Gam)=3$.
\end{proof}

\subsection*{$\Kinv(\Gam)=4$}

Algorithm~1 is deterministic once a labelling is fixed, and there are only
$3!\,4!=144$ labellings per orientation, so the value \eqref{eq:K} is a finite
computation on a graph with $8$ edges. Carried out over all of them, the number
of parts takes the values
$\{4,5\}$ on the orientation printed in \eqref{eq:gamma} and the constant value
$4$ on the orientation with $X$ and $Y$ exchanged. Hence
\[
  \Kinv(\Gam)=4
\]
in both readings, and with $\pmd(\Gam)=3$ this proves Theorem~\ref{thm:main}.
For instance the printed labelling itself gives the four parts
$\{x_2y_1,x_3y_2\}$, $\{x_1y_1,x_3y_3\}$, $\{x_1y_2,x_2y_3,x_3y_4\}$,
$\{x_2y_4\}$. The exhaustion is carried out by \texttt{verify.py} in exact
integer arithmetic.

\section{Verification}

The value $\pmd(\Gam)=3$ can be checked by hand from what is printed above: the
eight edges \eqref{eq:gamma}, the three parts $M_1,M_2,M_3$,
Lemma~\ref{lem:forest} and the bound \eqref{eq:delta}. The value
$\Kinv(\Gam)=4$ cannot. Its lower bound is the exhaustion of the $144$,
respectively $288$, labellings in \eqref{eq:K}, and we have no argument for it
that a reader can carry out without a machine.

The accompanying program \texttt{verify.py} (Python~3.9 or later, standard
library only, no external data file) reads the objects printed here as literals
and re-derives every quantity this paper states: the decoding of the graph6
string \texttt{FEhb?} and its agreement with \eqref{eq:gamma}; the
identification of $\Gam$ as $\Theta(2,2,4)$; the agreement of two structurally
independent implementations of \eqref{eq:pos}; the positivity of $M_1,M_2,M_3$
in their successive residuals and the value $\pmd(\Gam)=3$ by exhaustive
search; and $\Kinv(\Gam)=4$ over $144$ and over $288$ labellings. It prints one
line per check, closes with \texttt{VERDICT: ALL 80 CHECKS PASS} and exits $0$
only if every check passed.

The recorded run establishes more than this paper claims, and its labelling was
not brought back into line with the present text. Steps 1--6 are the checks that
bear on Theorem~\ref{thm:main}. Steps 7 and 8 --- a census of the bipartite
graphs of order at most $7$, the minimality and uniqueness of $\Gam$ among them,
and a second, cubic witness $M_5$ on ten vertices --- concern statements this
paper does not make, and their labels, and the closing \texttt{NOTE SCOPE}
paragraph, speak of values ``the paper prints'' that are printed in no version
of this note. Those labels are inherited from a longer manuscript and should be
read as claims about the program's own computation, not about anything asserted
here. Nothing in Sections~1--2 depends on them, and this note asserts none of
them.

Two things are taken from \cite{FGY} rather than checked here. The
characterisation \eqref{eq:pos} is a theorem of the source, used as the
definition of ``positive'' throughout; nothing here returns to the algebraic
definition. And Algorithm~1 has to be read: the program implements it both as
the source's pseudocode, statement for statement in the source's own variables,
and in the streamlined form used for the exhaustions, checks the two against
each other, and finds that both reproduce, part for part, the parts the source
prints for each of its two worked examples; but that the statements displayed in
Section~1 are the ones printed at lines $549$--$574$ of the e-print source of
\cite{FGY} is a comparison left to the reader.

\begin{thebibliography}{9}

\bibitem{FGY}
M.~Farrokhi D.~G., Sh.~Gharakhloo, and A.~A. Yazdan Pour,
\emph{Positive matching decompositions of graphs},
Discrete Appl. Math. \textbf{320} (2022), 311--323.
\href{https://doi.org/10.1016/j.dam.2022.05.012}{doi:10.1016/j.dam.2022.05.012};
\href{https://arxiv.org/abs/2110.12168}{arXiv:2110.12168}. Statement numbers
and line references are quoted from the e-print v2, the only full text we could
access.

\end{thebibliography}

\end{document}
